\section{An explicit family over unrestricted integers}\label{sec:integer}
\subsection{Fixed coefficients and base coordinates}
The six fixed coefficients are the coherent source-34 tuple
\[
 \beta=2^{1064921}-1,\quad K_0,\quad M_C,\quad M_F,
 \quad A_c=2^{1113},\quad B_c=2^{666104}.
\]
Here $M_F$ already includes the native shift by $\beta$. The complete
minimal little-endian numeral payloads and their binding file are supplied.
They are not unspecified parameters. The retained finite tile recipe and
source pin explain their mathematical origin; no choice based on halting
is involved. Their uncompressed total is 615,883 bytes.

Use the thirteen base coordinates
\[
 j,C,z,h,r,v,w,T,c,n,t,u_1,u_2\in\Z.
\]
Uppercase $C$ and lowercase $c$ are different coordinates. Set
\begin{align}
 J&=6j+1,&q&=\beta J+1,&W&=12w+4,\nonumber\\
 X&=qW,&Y&=qv,&E&=XY,\nonumber\\
 a&=(X+1)Y,&H&=4a+3,&\Delta&=(a+2)^2-1,\label{eq:front}\\
 F&=(K_0+W)C-z(q-1),\nonumber\\
 R&=[q(q-F)-C](q^2-1)+(M_C+qM_F)J,\nonumber\\
 k&=R+1+hE,&\eta&=c-Yk,&\zeta&=k-\eta,\label{eq:rank}\\
 I&=(2r+1)c^2,&D_f&=1+\Delta I^2,\nonumber\\
 D_\tau&=1+XY^2(XY^2+1)k^2,&D_m&=1+\Delta c^2,&b&=X+ac.
 \label{eq:radicands}
\end{align}
Define the determinant and positivity polynomials
\begin{align}
 G&=(A_cF)(B_cC)-(q-A_cF-B_cC)^2,\nonumber\\
 N&=G\eta\zeta-(t^2-G-\eta\zeta+1)^2,\nonumber\\
 S&=24N-21-(6u_1-1)^2-(6u_2-1)^2.\label{eq:guard}
\end{align}
The carrier scales are $M=2q^2$ and $U=MH$.
For the thirteen-variable equation put
\[
 Q=\Delta^2 I^2,\qquad V=D_f(cT-R)-c.
\]
For equations with a retained coordinate $f$, put instead
\begin{equation}
 Q=\Delta(f^2-1),\qquad V=f^2(cT-R)-c.\label{eq:retainedf}
\end{equation}
These agree when $f^2=D_f$. Finally set
\begin{equation}
 \delta=Q-1,\qquad C_y=QV^2-1,\qquad
 Z=U(n+V)+Mb+H(v-q^2).\label{eq:carrier}
\end{equation}

\subsection{A finite polynomial definition of the radical merger}
For $s$ radicands define the universal polynomial
\[
 \F_{2^s}(x;A_1,\ldots,A_s)=
 \prod_{\varepsilon\in\{-1,1\}^s}
 \left(x-\sum_{i=1}^s\varepsilon_i\sqrt{A_i}\right).
\]
Changing any radical sign permutes the factors, so this lies in
$\Z[x,A_1,\ldots,A_s]$. Each radicand has exponent at most $2^{s-1}$.
Write
\[
 \mathcal C_s(x;A_1,\ldots,A_{s-1};B,\delta)
 =\delta^{2^{s-1}}\F_{2^s}(x;A_1,\ldots,A_{s-1},B/\delta).
\]
The denominator is canceled symbolically. This is a finite polynomial
definition, including when $\delta=0$. The literal JSON circuits use
only addition, subtraction, and multiplication.

With $B_y=U^2C_y$, specialize
\begin{align*}
 N_5&=\mathcal C_5(Z;M^2D_m,9U^2D_f,81U^2D_\tau,729U^2S;B_y,\delta),\\
 N_4&=\mathcal C_4(Z;M^2D_m,9U^2D_\tau,81U^2S;B_y,\delta),\\
 N_3&=\mathcal C_3(Z;M^2D_m,9U^2D_\tau;B_y,\delta),\\
 N_2&=\bigl[\delta(Z^2-M^2D_m)-B_y\bigr]^2-4\delta M^2D_mB_y.
\end{align*}
The weights change after retaining the strong root: there is no reason
to preserve a vacant valuation position. This lowers the literal cost.
Put
\begin{align*}
 F_f&=f^2-\Delta I^2=1+(f^2-D_f),\\
 F_g&=g^2+22+(6u_1-1)^2+(6u_2-1)^2-24N=1+(g^2-S),\\
 F_\tau&=\tau^2-XY^2(XY^2+1)k^2=1+(\tau^2-D_\tau).
\end{align*}
The four equations are
\begin{align}
 P_{13}&=N_5=0,\label{eq:p13}\\
 P_{14}&=(1+N_4)F_f-1=0,\label{eq:p14}\\
 P_{15}&=(1+N_3)F_fF_g-1=0,\label{eq:p15}\\
 \boxed{P_{16}}&\boxed{=(1+N_2)F_fF_gF_\tau-1=0.}\label{eq:p16}
\end{align}
The added coordinates are none, $f$, $(f,g)$, and $(f,g,\tau)$,
respectively. All coordinates are unrestricted integers. The last line,
together with the polynomial abbreviations above, is the recommended
modest-degree equation in this report.

\subsection{Recovering retained equations by residues}
The affine formulas imply that $q$ is even, $q\equiv2\pmod3$, and
$W$ is a nonzero multiple of four with $W\equiv1\pmod3$. Thus $X$ is
divisible by eight, $Y$ and $a$ are even, and $\Delta\equiv3\pmod4$.
The integer $D_f$ is $0$ or $1\pmod4$, while $D_\tau\equiv1\pmod4$.
Lemma~\ref{lem:positive} gives $S\equiv1\pmod4$. Therefore
\[
 F_f\equiv f^2+I^2\in\{0,1,2\},\quad
 F_g\equiv g^2\in\{0,1\},\quad F_\tau\equiv\tau^2\in\{0,1\}\pmod4.
\]
None can equal $-1$.

Both $M$ and $U$ are divisible by eight. Modulo four every norm monomial
using a radicand vanishes, and
\begin{equation}
 N_s\equiv\delta^{2^{s-1}}Z^{2^s}\in\{0,1\}\pmod4\qquad(s\ge2).
 \label{eq:normmod4}
\end{equation}
There is no assumption that $\delta$ is odd; it need not be.
Lemma~\ref{lem:units} makes each product equation exactly its displayed
simultaneous system. In particular $f^2=D_f$ is recovered before using
\eqref{eq:retainedf} as the normalized strong Pell expression.

\subsection{Square extraction before positivity}
The order of the proof matters. Since $J\ne0$ and $\beta$ is large,
$|q|\ge16$ and $|X|\ge64$. If $v=0$, then $a=0$ and $\Delta=3$.
Otherwise $|Y|\ge16$ and $|a|=|X+1||Y|\ge1008$, so $\Delta>1$.
Always $H\ne0$ and $\Delta\ge3$.

After recovery of retained factors, if any, $Q=\Delta^2I^2$.
For $c=0$, $Q=0$ and $\delta=-1$; for $c\ne0$, $Q>1$ and
$\delta>0$. Thus the cleared norm can be interpreted as the product of
conjugate linear factors. Alternatively, the universal congruence
$\delta\equiv-1\pmod3$ already excludes a zero denominator.

Modulo three we have $X\equiv2$, $a=(X+1)Y\equiv0$, and
\[
 H\equiv\Delta\equiv0,\quad
 D_m\equiv D_f\equiv D_\tau\equiv S\equiv1,\quad
 Q\equiv0,\quad\delta\equiv C_y\equiv-1.
\]
For $D_\tau$, use $XY^2(XY^2+1)\equiv0\pmod3$ for either possible
residue of $Y^2$. Consequently $M$ is a unit at three and
$e=\vp(U)=\vp(H)\ge1$. In the five-root case the valuations of the
scaled radicands, in the order main, auxiliary, strong, first, guard, are
\[
 0,\quad2e,\quad2e+2,\quad2e+4,\quad2e+6.
\]
In the retained-strong family they are $0,2e,2e+2,2e+4$, omitting
any further retained roots. These are distinct even integers.
The norm zero gives a rational signed sum. Lemma~\ref{lem:valuation}
makes each unretained root rational, hence integer when its radicand is
integer. The guard acquires a root in this way or has a retained root.
Lemma~\ref{lem:positive} and the pair guard then give
\begin{equation}
 N>0,\qquad G>0,\qquad\eta\zeta>0.\label{eq:recoveredguard}
\end{equation}
If $c=0$, then $\eta\zeta=-Y(Y+1)k^2\le0$, a contradiction.
Thus $c\ne0$ and $Q>1$ at every solution.

\begin{lemma}[Adjacent-coefficient rational descent]\label{lem:descent}
If $Q>1$ and $V$ are integers and $y\in\Q$ satisfies
$QV^2-(Q-1)y^2=1$, then $y\in\Z$.
\end{lemma}
\begin{proof}
Write $y=A/d$ in lowest terms; then $d^2\mid Q-1$.
For $|V|=1$ the claim is immediate. After changing signs if needed, apply
\[
 V'=(2Q-1)V-2(Q-1)y,\qquad y'=(2Q-1)y-2QV.
\]
The first coordinate is integral, the norm is unchanged, and the reduced
denominator of the second remains $d$, since $2Q-1\equiv1\pmod d$.
For $V>1$, the norm gives $y>V$ and $(Q-1)y<QV$, whence
$-V<V'<V$; it also excludes $V'=0$. Descent on $|V|$ reaches one
without changing the denominator.
\end{proof}
The auxiliary radical has square $C_y/\delta$ and supplies the rational
$y$ in this lemma. It is therefore integral too.

\subsection{The two congruences and their inverse map}
The vanishing factor has the form
\begin{equation}
 Z=Md+UL,\qquad d^2=D_m,\quad L\in\Z,\label{eq:linearcarrier}
\end{equation}
where signs have been absorbed into the recovered roots. The integer
$L$ is their appropriate weighted sum. Because
$H=4(X+1)qv+3\equiv3\pmod q$, $H$ is odd, and $3\nmid q$,
\[
 \gcd(H,M)=\gcd(H,2q^2)=1.
\]
Modulo $M$, equation~\eqref{eq:linearcarrier} gives
\[
 v=q^2(2v_{\rm old}+1),\qquad
 v_{\rm old}=\frac{v-q^2}{2q^2}\in\Z.
\]
Thus $v\ne0$ and $Y=(2v_{\rm old}+1)q^3$, with exactly the required
odd quotient. Modulo $H$ it gives
$\gamma=(d-b)/H\in\Z$. Neither conclusion was assumed in the
valuation argument. Conversely choose
\begin{equation}
 n=\gamma+L-V-\frac{v-q^2}{M}\in\Z.\label{eq:choose-n}
\end{equation}
In the retained family $L=y+3\tau+9g$, omitting the terms whose roots
are retained; in the five-root family $L=y+3f+9\tau+27g$.

\subsection{The normalized source and its calibration}
The restored source has the same front registers and
\[
 f^2=D_f,\quad\tau^2=D_\tau,\quad QV^2-(Q-1)y^2=1,\quad
 (b+H\gamma)^2=D_m,\quad S=g^2,
\]
with $Y=(2v_{\rm old}+1)q^3$. The transport row holds identically
by the definition of $F$. Put $i=2r+1$ and restore older coordinates by
\begin{align*}
 I_{\rm old}&=\Delta i c^2,&r_{\rm old}&=(\Delta i-1)/2,\\
 o_{\rm old}&=f(cT-R),&j_{\rm old}&=D_fT-1-R\Delta i^2c^3.
\end{align*}
The first quotient is integral since $\Delta i$ is odd. Exact identities
give $I_{\rm old}^2=\Delta(f^2-1)$ and
$fo_{\rm old}-c=cj_{\rm old}-R=V$.
To restore the older ternary positivity row, use $N>0$ to choose
integers $b_1,b_2,b_3$ with
$N=1+b_1^2+b_2^2+b_3(b_3+1)$, by the three-square theorem.
These are reselected coordinates of the inherited source witness, not
additional variables in the new polynomial.

For completeness, the imported canonical native witness has
$q=2^e$ with $e$ odd, $R\equiv3\pmod4$, $X=2^R$, and an odd positive
quotient $Y/q^3$. Its native repunit satisfies $J\equiv1\pmod6$, and
$X/(4q)\equiv1\pmod3$. Thus the affine choices $J=6j+1$ and
$W=12w+4$ preserve canonical existence. These are properties of the
same complete source-34 coefficient tuple, not assumptions on arbitrary
integer assignments. Furthermore $c=\psi_A(R)$ with $A=a+2$ even
and $c$ odd. Choose
$I_*=\psi_A(Rc)$ and $f=\chi_A(Rc)$.
Pell composition, or expansion of $(D+c\sqrt\Delta)^c$, gives
$c^2\mid I_*$. The recurrence gives $I_*$ odd; hence $i=I_*/c^2$ is odd.
The retained squared-modulus auxiliary lemma, using the established
phase $R\equiv3\pmod4$, gives $V,y$ with
$V\equiv-c\pmod{f^2}$ and $V\equiv-R\pmod c$. Since
$f^2\equiv1\pmod c$, it supplies the integer
\[
 T=\frac{V+c+Rf^2}{cf^2}.
\]
The determinant and strip guards hold for that same canonical witness.
Choose $t$ by the nearest-square argument, the pentagonal representation
of $N$, and then $n$ by \eqref{eq:choose-n}.
This is the same calibrated source family as the retained root-fusion and
fixed-prime proposals, followed by the new algebraic finalizer.

\begin{theorem}[Integer family]\label{thm:integer}
Each polynomial \eqref{eq:p13}--\eqref{eq:p16}, with its stated fixed
tuple and domain, has a zero exactly when the normalized native source
has a witness. Composed with the inherited ordinary native/K215
calibration, absence of zeros is provably equivalent in ZFC to
$\Con(\mathrm{ZFC})$.
\end{theorem}
\begin{proof}
The product finalizer recovers retained equations; valuation separation
recovers other roots; descent recovers the auxiliary integer; the coprime
carrier recovers the cubic scale and main congruence. The displayed
restorations give the source. The preceding canonical construction proves
the reverse implication. Compose the named ordinary calibration and
negate. The imported compiler theorem is not replaced by numerical tests.
\end{proof}
